\section{Targets and reporting protocols}
\label{sec:methods}
\subsection{What does a selected score estimate?}
Let $j\in\{0,\ldots,K-1\}$ index a configuration fixed independently of the evaluation questions. A configuration includes the model, prompt, decoding, and scoring rule. For a question $X\sim P$, let $Y_j(X)\in\{0,1\}$ denote correctness and $\theta_j=\mathbb{E}_P Y_j(X)$ its accuracy. Scores for different configurations on the same question can be arbitrarily dependent. For the controlled and empirical-cache experiments below, search questions $S=(X_1,\ldots,X_n)$ are iid from the specified distribution, and exhaustive search returns
\begin{equation}
 \hat j=\arg\max_j\hat\theta_{Sj},\qquad
 \hat\theta_{Sj}=\frac1n\sum_{i=1}^nY_j(X_i),
\end{equation}
with independent uniform tie breaking. The reused report is $R=\hat\theta_{S\hat j}$, and its realized error is $R-\theta_{\hat j}$. Expected optimism averages this error over complete reruns of search. The selected target $\theta_{\hat j}$ is distinct from procedure quality $\mathbb{E}_S\theta_{\hat j}$ and the oracle $\max_j\theta_j$. A positive realized error is not guaranteed on every run.

For a fixed bank and iid bounded scores,
\begin{equation}
 0\leq\mathbb{E}[R-\theta_{\hat j}]
 \leq\sqrt{\frac{\log K}{2n}}.
 \label{eq:optimism}
\end{equation}
The lower bound follows from $\mathbb{E}\max_j\hat\theta_{Sj}\geq\max_j\theta_j\geq\mathbb{E}\theta_{\hat j}$. A bounded moment-generating-function argument gives the upper bound (Appendix~\ref{app:guarantees}). This expectation bound is often loose, particularly with strongly shared errors. It does not imply that optimization cannot improve true quality.

\subsection{Reporting alternatives}
\paragraph{Reuse.}
Report the winning search accuracy with its ordinary two-sided 95\% Clopper--Pearson (CP) interval. CP has a fixed-candidate binomial guarantee; selecting the candidate using the same outcomes does not preserve that guarantee.

\paragraph{Simultaneous fixed-bank protection.}
Construct CP intervals for every candidate at marginal error $0.05/K$, then report the interval for $\hat j$. A union bound gives simultaneous coverage at least 95\%, hence at least 95\% coverage for the selected member, without independence across candidates. This changes the interval, not the reused point estimate. The guarantee requires an outcome-independent family; unrestricted feedback-based creation of new candidates is a different problem.

\paragraph{Independent audit.}
Freeze the selected configuration before evaluating it on $h$ iid fresh questions. The audit score is conditionally unbiased for $\theta_{\hat j}$, with variance $\theta_{\hat j}(1-\theta_{\hat j})/h$ and ordinary CP coverage at least 95\%. Further selection after viewing audit outcomes would change the argument. We compare an additional audit of size $n$ after full search with two split protocols. A question split reserves $h=n-\lfloor n/2\rfloor$ questions and searches on $s=\lfloor n/2\rfloor$, preserving $n$ total question draws. A call split takes
\begin{equation}
 s=\left\lfloor\frac{nK}{K+1}\right\rfloor,\qquad h=nK-sK,
\end{equation}
so $sK+h=nK$ preserves total scoring calls. The latter can use more question draws. Both splits can select a different configuration from full search; their reports target their own winners.

\subsection{How reliability is evaluated}
We rerun selection in each repetition and measure signed error, mean squared error (MSE), interval coverage and width, and selected true quality. All intervals concern the selected configuration's accuracy, with marginal coverage over the random selected identity. We do not claim identity-conditional coverage from a marginal result.

For a concrete decision endpoint, candidate 0 is a frozen default. Declare superiority when a nondefault winner's ordinary 95\% CP lower endpoint exceeds the default's corresponding upper endpoint. A false claim has $\theta_{\hat j}\leq\theta_0$. The two relevant one-sided errors sum to at most 0.05 for a frozen pair; that argument need not hold after choosing the pair by search. We also evaluate the simultaneous-interval version. We report unconditional false-superiority frequency, not the fraction of claims that are false, and do not identify this interval-nonoverlap rule with a paired test.


\subsection{Finite-benchmark and empirical-cache targets}
The primary real-benchmark protocol uses sampling without replacement rather than the iid experiment above. A prespecified runtime fallback retains four candidates, 256 search questions and 512 disjoint audit questions from the pinned MMLU-Pro benchmark's 12,032 test items. Conditional on the retained search reservoir $S_C$, the target for a chosen configuration is its question-weighted accuracy on the remaining 11,776 items, $\theta_{\rm rem}(\hat j;S_C)$. Smaller search subsets retain this same complement target. The audit score carries sampling uncertainty and is not known truth.

For each fixed design, exact hypergeometric inversion supplies audit intervals and transformed simultaneous-search intervals for this target. However, runtime can depend on item properties related to correctness. Outcome-blind choice between the original and fallback designs therefore does not by itself preserve fixed-design guarantees conditional on that choice. The supplied analysis allocates error $0.05/2$ to each possible design, with a further bank adjustment where needed. A union bound protects the runtime-chosen interval at least 95\% unconditionally over partition and tier choice, separately for each endpoint family. We claim neither coverage nor audit unbiasedness conditional on the runtime trigger, nor joint protection of all reported endpoints. Appendix~\ref{app:real-protocol} gives the sampling argument, transformation, scope history and costs.

A separate calibration experiment samples whole paired rows iid with replacement from the 768-row retained response cache, rerunning selection and reporting. Candidate means in that empirical distribution are known, permitting coverage and decision error to be measured against its conditional truth with the original 95\% iid protocols. This does not establish calibration on the unobserved remainder of the benchmark. Reusing one audit pool likewise cannot remove its sampling uncertainty by increasing the number of Monte Carlo repetitions.
